
on
définit un \(\lambda\)-homomorphisme : \(B\to K\) prolongeant l'augmentation de \(A\) en envoyant \(X\) sur \(0\), ce qui munit \(B\) d'une structure de \(K\)-\(\lambda\)-algèbre augmentés. On se propose de déterminer sa filtration (3.10) en fonction de celle de \(A\).

On gradue \(B\) par le poids par rapport aux \(\gamma^k X\), \(\gamma^k X\) étant affecté du poids \(k\)~; on désigne par \(M_k\) l'ensemble des monômes de poids \(k\) par rapport aux indéterminées \(\gamma^i X\).

\textbf{Lemme 4.11.} \uline{Soit} \(a\in A\)~; \uline{le coefficient d'un monôme de poids} \(n\) \uline{par rapport} \uline{aux} \(\gamma^j X\) \uline{dans} \(\gamma^k(aX)=P'_k(\gamma^i(a),\gamma^j X)\) \uline{est de filtration au moins} \(k-n\) (\uline{dans} \(A\))

Le polynôme \(P'_k\) est sans composantes isobares de poids \(<k\) par rapport aux \(\gamma^i(a)\), \(\gamma^j X\) (3.5) donc le lemme est vrai si \(a\) est d'augmentation nulle.

Si \(\epsilon\) désigne l'homomorphisme d'augmentation : \(A\to K\), on peut écrire dans le cas général :
\[
  \gamma_t(aX)=\gamma_t((a-\epsilon(a))X)\mathbin{.}\gamma_t(\epsilon(a)X).
\]
Or \(\gamma_t(\epsilon(a)X)=(\gamma_t(X))^{\epsilon(a)}=\sum_{n=0}^{\infty}R_{n,a}((\gamma^j X))\mathbin{.}t^n\), les \(R_{n,a}\) étant isobares de poids \(n\) par rapport aux \(\gamma^j X\) (3.8)~; on en déduit le lemme en faisant le produit.

\textbf{Proposition 4.12.} \(\Fil^k B=\Fil^k A+\Fil^{k-1}A\mathbin{.}M_1+\ldots+\Fil^0 A\mathbin{.}M_k+B\mathbin{.}M_{k+1}\).

Il est clair que le second idéal est contenu dans le premier~; il y a donc simplement à montrer la réciproque.

Comme la filtration définie par la famille des idéaux du second membre pour \(k\) variable est une filtration d'anneau, il suffit de montrer la réciproque pour un élément du type \(\gamma^k(b)\), avec \(\epsilon(b)=0\), ces éléments
